% Seção: Experimentos e resultados  (orçamento: 2.5 página(s))
%
% Narrativa combinada com o autor (13/10/2026): Ferramentas mantêm o acerto nas duas faixas de preço; Consulta e
% Contexto longo dependem de modelos de preço alto, com raciocínio, e pagam em tokens de saída e latência; tool calling
% efetivo é condição. Os Resultados descrevem; a Discussão interpreta. Sem McNemar no texto por ora (a tabela
% generated/tables/mcnemar_models.tex continua gerada). Sem preço em dólar (os preços mudam): só tokens e latência.

\section{Experimentos e resultados}\label{sec:experimentos}

\subsection{Descrição dos dados}

Os dados são os eventos da final da Copa do Mundo de 2022: \nEvents{} eventos no JSON bruto, que viraram
\nActions{} ações na tabela de ações. As amostras do experimento são as perguntas: \nQuestions{} perguntas em cinco
grupos (Tabela~\ref{tab:perguntas}), das quais \nAnswerable{} têm resposta nos dados e \nControl{} formam o grupo de
controle. Cada pergunta foi feita a cada uma das \nArms{} abordagens, num total de \nRuns{} execuções por modelo.

\vspace{2mm}
\noindent\begin{minipage}{\columnwidth}
\centering\small\renewcommand{\arraystretch}{1.15}
\begin{tabular}{>{\raggedright\arraybackslash}p{0.28\columnwidth}>{\raggedright\arraybackslash}p{0.60\columnwidth}} % lint: ignore
\toprule
Grupo & Exemplo (resposta esperada) \\
\midrule
Busca e contagem (\nGroupSearch{}) & Quem recebeu o primeiro cartão amarelo da final? (jogador) \\
Relações na rede de passes (\nGroupRelations{}) & Qual dupla da França mais trocou passes entre si na final? (conjunto de jogadores) \\
Jogadas e sequências (\nGroupPlays{}) & Qual jogador participou de mais jogadas da Argentina que terminaram em finalização? (jogador) \\
Antes e depois de um momento (\nGroupMoments{}) & Depois que Dembélé foi substituído no primeiro tempo, qual companheiro mais passou a bola para Mbappé? (jogador) \\
Controle (\nGroupControl{}) & Qual foi a distância total percorrida por Rodrigo De Paul na final? (ausência de dados) \\
\bottomrule
\end{tabular}
\captionof{table}{Os cinco grupos de perguntas, com o número de perguntas de cada um e um exemplo.}\label{tab:perguntas}
\end{minipage}\par\smallskip

\subsection{\textit{Setup} experimental}

Foram avaliados \nModels{} modelos de linguagem, em duas faixas de preço de tabela (outubro de 2026). Na faixa
baixa, de centavos de dólar por milhão de tokens, ficaram \modelHaiku{} e \modelQwen{}; na alta, de alguns dólares
por milhão, \modelSonnet{} e \modelKimi{}. Em cada faixa há um modelo de código fechado e um de pesos abertos. Cada
modelo respondeu a todas as perguntas em todas as abordagens, com os mesmos prompts. Os modelos de preço baixo rodaram
sem raciocínio estendido e com temperatura zero; os de preço alto, com raciocínio estendido no nível padrão e na
temperatura padrão do provedor. A resposta de cada chamada ao modelo, raciocínio incluído, foi limitada a
\nMaxOutTokens{} tokens. Na abordagem \armEventsInPrompt{} do \modelSonnet{}, o raciocínio no nível padrão esgotou
esse limite em parte das perguntas, que ficaram sem resposta; nessa abordagem, o nível de raciocínio foi reduzido, e
as perguntas que ainda esgotaram o limite foram refeitas com o dobro dele.

A abordagem \armVector{} usou embeddings do modelo \texttt{\embedModel}. Nas abordagens com agente, cada pergunta
admitiu no máximo \nMaxCalls{} chamadas de ferramenta. A rodada de cada modelo teve \nRepeats{} repetição por
pergunta.

% TODO autor: REVISAR este parágrafo (movido da Metodologia; ainda não decidido se entra no artigo)
\begingroup\color{red}\textbf{[REVISAR]} O acerto de uma abordagem é a proporção de perguntas respondidas corretamente. Como as perguntas são poucas e cada uma
foi executada uma vez, duas análises acompanham essa proporção, ambas calculadas depois da rodada, a partir do
arquivo de resultados. A primeira é o intervalo de confiança de Wilson~\cite{wilson1927probable}, de \ciLevel{}
para cada acerto; diferentemente do intervalo normal, ele não ultrapassa zero nem cem por cento, o que importa para
as abordagens em zero por cento. A segunda compara pares de abordagens que responderam às mesmas perguntas: o teste exato
de McNemar~\cite{mcnemar1947note} considera apenas as perguntas em que uma abordagem acertou e a outra errou, e é o que
Dietterich~\cite{dietterich1998approximate} indica quando o algoritmo é executado uma só vez. Foram comparados
\armTextToCypher{} e \armEventsInPrompt{}, \armGraphTools{} e \armTextToCypher{}, e \armGraphTools{} e
\armEventsInPrompt{}. Como são três testes, os valores $p$ recebem o ajuste de Holm~\cite{holm1979simple}, que
mantém no nível escolhido a probabilidade de algum falso positivo no conjunto. A rodada tem \nRepeats{} repetição
por pergunta; as conclusões são lidas como indícios, e não como demonstração.\par\endgroup

% TODO autor: REVISAR (confirmar se o repositório é público e o link; o template pede aqui as etapas para reproduzir)
\begingroup\color{red}\textbf{[REVISAR]} O código, as perguntas, o gabarito, os resultados de cada execução e uma
página por rodada estão em \url{https://github.com/joaowinderfeldbussolotto/football-tactics-knowledge-graph}. As
bibliotecas usadas foram socceraction~\cite{socceraction}, Neo4j Graph Data Science~\cite{neo4j_gds},
networkx~\cite{networkx}, PydanticAI~\cite{pydanticai} e Langfuse~\cite{langfuse}, nas versões indicadas nas
referências.\par\endgroup

\subsection{Análise dos resultados}

\subsubsection{Acerto por abordagem e modelo}

A Tabela~\ref{tab:acerto} mostra o acerto das \nArms{} abordagens nos \nModels{} modelos, e a
Figura~\ref{fig:modelos} destaca as três que leem os dados da partida. \armGraphTools{} acertou de
\accRangeGraphTools{} das perguntas com resposta, nas duas faixas de preço. \armTextToCypher{} acertou de
\accRangeLowTextToCypher{} nos modelos de preço baixo e de \accRangeHighTextToCypher{} nos de preço alto;
\armEventsInPrompt{}, de \accRangeLowEventsInPrompt{} e de \accRangeHighEventsInPrompt{}. Nos modelos de preço
baixo, \armGraphTools{} ficou à frente das outras duas, com intervalos que não se sobrepõem aos de
\armEventsInPrompt{}; nos de preço alto, os intervalos das três se sobrepõem. As outras abordagens ficaram
abaixo em todos os modelos: \armVector{} acertou de \accRangeVector{}, \armStatsInPrompt{} de
\accRangeStatsInPrompt{} e \armNoContext{}, \accRangeNoContext{}.

\vspace{2mm}
\noindent\begin{minipage}{\columnwidth}
\centering\footnotesize\renewcommand{\arraystretch}{1.15}\setlength{\tabcolsep}{4pt}
\gentable{models}
\captionof{table}{Acerto, em porcentagem, nas perguntas com resposta. Os intervalos de confiança das três
abordagens que leem os dados estão na Figura~\ref{fig:modelos}.}\label{tab:acerto}
\end{minipage}\par\smallskip

\vspace{2mm}
\noindent\begin{minipage}{\columnwidth}
\centering
\gentable{fig_models}
\captionof{figure}{Acerto das três abordagens que leem os dados nas perguntas com resposta, por modelo, com o
intervalo de Wilson de \ciLevel{}. A área cinza marca os modelos de preço alto.}\label{fig:modelos}
\end{minipage}\par\smallskip

\subsubsection{Acerto por grupo de perguntas}

A Tabela~\ref{tab:grupos} mostra os acertos das três abordagens em cada grupo de perguntas
(Seção~\ref{sec:perguntas}). Nos modelos de preço baixo, a distância entre \armGraphTools{} e as outras duas
abordagens aparece sobretudo nos grupos de jogadas e de relações na rede de passes. Nos modelos de preço alto,
\armGraphTools{} acertou menos que \armTextToCypher{} no grupo antes e depois de um momento do jogo.

\vspace{2mm}
\noindent\begin{minipage}{\columnwidth}
\centering\footnotesize\renewcommand{\arraystretch}{1.15}\setlength{\tabcolsep}{4pt}
\gentable{groups_models}
\captionof{table}{Acertos por grupo de perguntas, em cada célula na ordem \armShortEventsInPrompt{},
\armShortTextToCypher{} e \armShortGraphTools{}. Busca: busca e contagem; Relações: relações na rede de passes;
Jogadas: jogadas e sequências; Momento: antes e depois de um momento do jogo.}\label{tab:grupos}
\end{minipage}\par\smallskip

\subsubsection{Erros e custo}

A Tabela~\ref{tab:erros} mostra, para as três abordagens que leem os dados, a proporção de alucinações, de
abstenções e de erros de formato (Seção~\ref{sec:conferencia}) e, por execução, a média de tokens de entrada e de
saída e a latência. Nos modelos de preço alto, \armEventsInPrompt{} produziu em média de
\outRangeHighEventsInPrompt{} tokens de saída por pergunta, quase todos de raciocínio, contra
\outRangeLowEventsInPrompt{} nos de preço baixo, e levou de \latRangeHighEventsInPrompt{} segundos. Os erros de
formato vieram de limites: no \modelQwen{}, o agente esgotou as chamadas de ferramenta em \armTextToCypher{}
(\fmtQwenTextToCypher{}) e em \armGraphTools{} (\fmtQwenGraphTools{}); no \modelKimi{}, em \armGraphTools{}
(\fmtKimiGraphTools{}), o modelo esgotou o limite de saída ou as tentativas de resposta; no \modelSonnet{}, em
\armEventsInPrompt{} (\fmtSonnetEventsInPrompt{}), esgotou o limite de saída mesmo com o dobro dele. Fora da tabela, \armNoContext{} e \armStatsInPrompt{} se abstiveram em
\abstRangeNoContext{} e \abstRangeStatsInPrompt{} das perguntas, e \armVector{}, em \abstRangeVector{}.

\vspace{2mm}
\noindent\begin{minipage}{\columnwidth}
\centering\footnotesize\renewcommand{\arraystretch}{1.15}\setlength{\tabcolsep}{3pt}
\gentable{errors_models}
\captionof{table}{Perfil de erro sobre todas as perguntas e médias por execução das três abordagens que leem os
dados, por modelo. Form.: erro de formato; entrada e saída em milhares de tokens; latência em segundos.}\label{tab:erros}
\end{minipage}\par\smallskip
